\documentclass[11pt]{article}
\usepackage[a4paper,margin=24mm]{geometry}
\usepackage{amsmath,amssymb,hyperref}
\title{The stated ceiling rank bound is not sharp at two constraints}
\author{ziangni-sys}
\date{4 October 2026}
\begin{document}
\maketitle
We disprove the pointwise sharpness of the ceiling formula in conjecture
00000008822. The displayed rank inequality and the asserted attainment
of the ceiling bound are incompatible at $m=2$. We do not use the
separate assertions about minimal-rank solutions or Frank--Wolfe.

For real symmetric semidefinite programs, write the actual feasible set as
\[
F(A,b)=\{X\succeq0:\operatorname{tr}(A_iX)=b_i,
\ i=1,\ldots,m\}.
\]
The conjecture asserts that the rank $r$ of any extreme point satisfies
\[
\frac{r(r+1)}2\le m,
\]
and claims that the bound
$c(m)=\lceil(\sqrt{8m+1}-1)/2\rceil$ is tight: an extreme point can
have rank exactly $c(m)$.

At $m=2$, we have $3<\sqrt{17}<5$, hence
\[
1<\frac{\sqrt{17}-1}{2}<2,\qquad c(2)=2.
\]
But an extreme point of that rank would, by the simultaneously asserted
rank inequality, satisfy
\[
3=\frac{2(2+1)}2\le2,
\]
a contradiction. More generally, for integer $r\ge0$ the displayed
inequality at $m=2$ forces $r\le1$. This argument covers every dimension,
all choices of the two trace constraints and every actual extreme point.
Thus no such instance can attain the stated ceiling value under the
conjecture's own rank law.

The correct integer upper bound obtained from the displayed inequality
uses the \emph{floor}, not the ceiling:
\[
r\le\left\lfloor\frac{\sqrt{8m+1}-1}{2}\right\rfloor.
\]
The ceiling is a weaker bound; calling it pointwise sharp is the error.
For triangular values $m=r(r+1)/2$, floor and ceiling coincide. We do
not deny that examples can attain the formula at those special values;
the counterexample concerns its asserted sharpness as a bound in $m$.

A concrete feasible example at $m=2$ is the scalar SDP with
$A_1=A_2=[1]$ and $b_1=b_2=1$. Its feasible set is the singleton
$\{[1]\}$, whose sole point is genuinely extreme and has rank $1$.
This illustrates the attainable rank without replacing the universal
contradiction by an example-specific assumption.

\paragraph{Formal correspondence.}
Lean defines $F(A,b)$ using the actual \texttt{Matrix.PosSemidef},
matrix multiplication and trace. It uses the actual
\texttt{Set.extremePoints} and \texttt{Matrix.rank}, quantifies over
all finite matrix dimensions and all two-constraint instances, computes
$c(2)=2$, and disproves the conjunction of the displayed universal rank
law with ceiling attainment. It also checks the explicit singleton SDP
and its rank-one extreme point. The inequality itself is a hypothesis
of the source conjunction being refuted; the Lean proof does not pretend
to reprove Pataki's theorem. Lean 4.19.0 and Mathlib commit
\texttt{c44e0c8ee63ca166450922a373c7409c5d26b00b} are pinned.
\end{document}
